%  =============================================================================
%  Datei-Hinweise & Konfiguration: siehe config.md
%  -> Abschnitt "anhang/evaluationsbelege.tex"
%
%  Anhang "Evaluationsbelege": reproduzierbare Belege der in Kapitel 3.4
%  genannten Kennzahlen. Druckaufbereitung der Protokolle analog zu
%  anhang/quelltexte.tex (ASCII-Ersetzungstabelle, footnotesize).
%  =============================================================================

\chapter{Evaluationsbelege}
\label{sec:anhang-evaluationsbelege}

Dieser Anhang dokumentiert die Datengrundlagen der Evaluation
(\cref{sec:umsetzung-evaluation}). Sämtliche Protokolle wurden unverändert von
den Produktionssystemen übernommen und lediglich typografisch für den Druck
aufbereitet; Umfangreiche Rohdaten (Tagesbars, Rohticks) sind durch
Herkunftsangaben referenziert, statt abgedruckt.

\section{Backtest-Protokoll (B0/B1/B2)}
\label{sec:anhang-backtest}

Vollständige Ausgabe von \texttt{strategy.py backtest} für die Kapitaleinsätze
2.000\,€ und 1.000\,€ (Datenstand 22.07.2026, 2.520 Handelstage; erzeugt am
23.07.2026 auf der agent-vm). Die in \cref{sec:umsetzung-evaluation-ergebnisse}
genannten Kennzahlen entstammen unverändert diesem Protokoll.

\lstinputlisting[basicstyle=\ttfamily\footnotesize,breaklines=true]{anhang/belege/backtest_protokoll_20260723.txt}

\section{Walk-Forward-Kalibrierung der statistischen Prüfinstanz}
\label{sec:anhang-walkforward}

Kalibrierungstabelle der Walk-Forward-Prüfung (\texttt{ml\_walkforward.csv};
23.803 bewertete Prognosetage): je Wahrscheinlichkeits-Bin die Anzahl der
Prognosen, die mittlere prognostizierte Wahrscheinlichkeit und die beobachtete
Trefferrate.

\begin{table}[htbp]
  \centering
  \begin{tabular}{r r r r}
    \hline
    Bin & $n$ & $\bar{p}$ prognostiziert & beobachtete Rate \\
    \hline
    0{,}1 & 2      & 0{,}184 & 0{,}000 \\
    0{,}2 & 7      & 0{,}262 & 0{,}429 \\
    0{,}3 & 127    & 0{,}364 & 0{,}504 \\
    0{,}4 & 4.142  & 0{,}483 & 0{,}501 \\
    0{,}5 & 19.434 & 0{,}524 & 0{,}525 \\
    0{,}6 & 91     & 0{,}624 & 0{,}549 \\
    \hline
  \end{tabular}
  \caption[Walk-Forward-Kalibrierung der statistischen Prüfinstanz]{%
    Walk-Forward-Kalibrierung der statistischen Prüfinstanz
    (Quelle: \texttt{ml\_walkforward.csv}, Stand 22.07.2026).\quelleeigen}
  \label{tab:anhang-walkforward}
\end{table}

\section{Ex-Ante-Transaktionskosten (Bankmessung)}
\label{sec:anhang-kosten}

Rohdaten der Kostenmessung gegen das Produktivsystem der Bank
(\texttt{cost\_probe.py}, Ex-Ante-Kostenausweis der Ordervalidierung, Status
201 = bankseitig validiert; Handelsplatz Live Trading/Direkthandel):

\begin{table}[htbp]
  \centering
  \begin{tabular}{r r r r r r r}
    \hline
    Volumen & Stück & Provision & Produktkosten & Kauf gesamt & Round-Trip & RT-Quote \\
    \hline
    150\,€ & 10 & 3{,}90\,€ & 0{,}14\,€ & 4{,}04\,€ & 7{,}80\,€ & 5{,}20\,\% \\
    495\,€ & 33 & 3{,}90\,€ & 0{,}47\,€ & 4{,}37\,€ & 7{,}80\,€ & 1{,}58\,\% \\
    990\,€ & 66 & 3{,}90\,€ & 0{,}95\,€ & 4{,}85\,€ & 7{,}80\,€ & 0{,}79\,\% \\
    \hline
  \end{tabular}
  \caption[Ex-Ante-Transaktionskosten der Bankmessung]{%
    Ex-Ante-Transaktionskosten der Bankmessung
    (Quelle: \texttt{kostentabelle.csv}, erhoben 19.07.2026).\quelleeigen}
  \label{tab:anhang-kosten}
\end{table}

\section{Herkunftsnachweis der Marktdatenbasis}
\label{sec:anhang-provenienz}

Herkunfts- und Integritätsnachweis der produktiven Marktdatenbank
(\texttt{marketdata.sqlite}, collector-vm; Abfragen vom 23.07.2026). Die
Rohtick-Tabelle enthält konstruktionsbedingt nur die jüngsten Handelstage, da
ältere Rohticks nach zehn Tagen zu Bars verdichtet und beschnitten werden
(vgl. \cref{sec:umsetzung-implementierung}).

\begin{table}[htbp]
  \centering
  \begin{tabular}{l l}
    \hline
    Merkmal & Wert \\
    \hline
    Tagesbars \texttt{bars\_1d} (S\&P~500) & 30.12.1927 -- 22.07.2026, 24.755 Zeilen \\
    Rohticks \texttt{ticks\_raw} (rollierend) & 79.945 Zeilen \\
    Integrität (\texttt{md5}) & \texttt{24771c148ec062b2c2dfa790b9d22975} \\
    Replikation (Litestream) & \texttt{gs://fom-ki-marketdata-backup/marketdata/} \\
    \hline
  \end{tabular}
  \caption[Herkunftsnachweis der Marktdatenbasis]{Herkunftsnachweis der
    Marktdatenbasis (collector-vm, Stand 23.07.2026).\quelleeigen}
  \label{tab:anhang-provenienz}
\end{table}

Die jüngsten Zeilen der Zweitquellenprüfung (\texttt{source\_crosscheck};
Spalte \texttt{stooq\_close} enthält aus historischen Gründen den
FRED-Schlusskurs) belegen die laufende Übereinstimmung der Datenquellen:

\begin{table}[htbp]
  \centering
  \begin{tabular}{l r r r r}
    \hline
    Datum & Primärquelle & FRED & Divergenz & Suspekt \\
    \hline
    21.07.2026 & 7.509{,}20 & 7.509{,}20 & 0{,}0\,\% & 0 \\
    20.07.2026 & 7.443{,}28 & 7.443{,}28 & 0{,}0\,\% & 0 \\
    17.07.2026 & 7.457{,}69 & 7.457{,}69 & 0{,}0\,\% & 0 \\
    16.07.2026 & 7.533{,}77 & 7.533{,}77 & 0{,}0\,\% & 0 \\
    15.07.2026 & 7.572{,}40 & 7.572{,}40 & 0{,}0\,\% & 0 \\
    14.07.2026 & 7.543{,}59 & 7.543{,}59 & 0{,}0\,\% & 0 \\
    13.07.2026 & 7.515{,}34 & 7.515{,}34 & 0{,}0\,\% & 0 \\
    10.07.2026 & 7.575{,}39 & 7.575{,}39 & 0{,}0\,\% & 0 \\
    09.07.2026 & 7.543{,}64 & 7.543{,}64 & 0{,}0\,\% & 0 \\
    \hline
  \end{tabular}
  \caption[Zweitquellenprüfung der Tagesschlusskurse]{Jüngste Zeilen der
    Zweitquellenprüfung gegen \glsxtrshort{fred}
    (Stand 23.07.2026).\quelleeigen}
  \label{tab:anhang-crosscheck}
\end{table}

% ---------------------------------------------------------------------------
% PLATZHALTER (wird nach Go-live 27.07. befüllt):
%
% \section{Belegkette der ersten Echtgeld-Order}
%   -> Tagesreport (redigiert), Journal-Auszug,
%      Ordervalidierung Status 200/201 (BELEGKONZEPT, Abschnitt B3).
% ---------------------------------------------------------------------------
